\documentclass[11pt]{article}
\usepackage[margin=1in]{geometry}
\usepackage[T1]{fontenc}
\usepackage{lmodern,amsmath,amssymb,amsthm,booktabs,microtype}
\usepackage[colorlinks=true,allcolors=blue]{hyperref}
\hypersetup{pdftitle={Capacity-loss timing and optimal inventory targets: a complete three-period classification}}
\newtheorem{theorem}{Theorem}[section]
\newtheorem{proposition}[theorem]{Proposition}
\newtheorem{lemma}[theorem]{Lemma}
\newtheorem{corollary}[theorem]{Corollary}
\theoremstyle{remark}\newtheorem{remark}[theorem]{Remark}
\newcommand{\EE}{\mathbb E}
\newcommand{\argmin}{\mathop{\rm arg\,min}}
\newcommand{\Ecase}{\mathsf E}
\newcommand{\Lcase}{\mathsf L}
\title{Capacity-loss timing and optimal inventory targets:\\a complete three-period classification}
\author{}
\date{}
\makeatletter
\renewcommand{\maketitle}{\begin{center}{\Large\bfseries\@title\par}\end{center}\vspace{0.5em}}
\makeatother
\begin{document}
\maketitle
\begin{abstract}
An earlier loss of production capacity need not increase the current
base-stock target more than an equally large later loss, contrary to a
published advance-capacity-information conjecture. We give a complete
classification for a three-period backlogged inventory model with
independent identically distributed two-point demand. For every
nondegenerate demand probability and every positive holding-to-backlog
cost specification, six explicit cost thresholds determine both target
sets, including all intervals of optimal real orders at equality.
Two open cost regions exhibit the reversed temporal ordering and two
exhibit the opposite ordering. Both signs attain the universal
one-loss target-spread bound. A low-backlog example also contradicts the
claim that both targets coincide when holding cost exceeds backlog cost.
The reversals persist under explicit continuous-density demand
perturbations. In a separately specified loss-location information
experiment, we compute exact value of information and minimax regret.
The classification follows from symbolic Bellman slopes; independent
exact and scenario-tree implementations accompany the analysis.
\end{abstract}
\noindent\textbf{Keywords:} inventory control; advance capacity information;
base-stock policy; comparative statics; value of information.\\
\textbf{2020 Mathematics Subject Classification:} 90B05, 90C39.

\section{Introduction}
Advance capacity information changes inventory decisions by revealing
future replenishment constraints. This raises two different comparative
statics questions: how much a target can change after a capacity loss,
and how the change depends on the time of that loss. A coordinate
sensitivity bound does not by itself order the effects of losses in
different periods.

Jak\v{s}i\v{c} et al.\ propose that an earlier equal capacity reduction
produces at least as high a current target
\cite[Conjecture 1, p.~360]{Journal}. The immediately preceding discussion
also asserts that the two targets coincide when $h>b$. We disprove both
assertions within the published model. The temporal comparison was
Theorem 4(1) in the 2008 working paper \cite{WP} and later appears as
Conjecture 2.1 in the 2016 thesis \cite{Thesis}. The final journal's
Theorem 4 concerns coordinatewise capacity monotonicity; our examples
do not contradict it.

We resolve the temporal comparison completely in one finite-horizon
family. For each demand probability $p\in(0,1)$, six explicit thresholds
partition \emph{all} positive cost ratios. Reversals occur both below and
above the myopic cost threshold. The examples attain both signs of the
sharp universal target-spread bound and remain valid for small
continuous-density demand perturbations. The result concerns target
decisions; it does not rank total optimized costs or welfare losses.

Other supply-information models have different primitives. Altu\u{g} and
Muharremo\u{g}lu \cite{Altug} model evolving capacity forecasts, whereas
Atasoy, G\"ull\"u and Tan \cite{Atasoy} consider deterministic
nonstationary demand with all-or-nothing availability. A recent model
allows capacity to carry over and regenerate \cite{Senga}; our unused
period capacity expires. These distinctions matter when transferring
policy comparisons. The present classification does not assert that
base-stock optimality or the general value of information is new.

\section{Model and complete classification}\label{sec:model}
There are three periods, zero lead time, nonnegative real orders,
full backlogging, no ordering cost, no terminal salvage and no discount.
The end-of-period loss is
\[
 \ell(x)=h x_++b(-x)_+,\qquad h,b>0.
\]
Demands are iid with $D_t=2\eta B_t$,
$\Pr(B_t=1)=p\in(0,1)$, and $\eta>0$.
An order placed in period $t$ is bounded by known capacity $q_t$.
For $V_4=0$, define
\begin{equation}\label{eq:bellman}
 J_t(y)=\EE[\ell(y-D_t)+V_{t+1}(y-D_t)],\qquad
 V_t(x)=\min_{0\le z\le q_t}J_t(x+z).
\end{equation}
The target set is $\argmin_yJ_t(y)$; the base-stock selector is its
smallest member. Targets are unconstrained post-order levels, even if
current capacity prevents reaching them. Starting inventory zero and
$q_1\ge4\eta$ make every target in this paper feasible.

Starting from future capacities $(\eta,2\eta)$, compare the early-loss
and late-loss vectors
\[
 q^{\Ecase}=(0,2\eta),\qquad q^{\Lcase}=(\eta,\eta).
\]
Both remove $\eta$ units and leave the same total future capacity.
They are compared as known future vectors. To embed all three vectors in
one independent capacity process, let $Q_2$ have support $\{0,\eta\}$
and $Q_3$ have support $\{\eta,2\eta\}$, with positive probabilities
and independent of demands. Each vector is then a possible fully revealed
ACI state. A two-point prior on \emph{the two loss vectors alone} is a
different experiment, specified in
Section~\ref{sec:information}.

Set $r=b/h$, $u=1-p$, $c=3-3p+p^2$ and $r_0=u/p$. Define
\begin{align}\label{eq:thresholds}
 A&=\frac{uc}{p(5-3p+p^2)},&
 B&=\frac{uc}{p(4-2p+p^2)},&
 C&=\frac{u(3-p)}{3p},\notag\\
 L&=\frac{3-p-p^2}{2p^2},&
 E&=\frac{3-2p+p^2-p^3}{p^2(1+p)},&
 U&=\frac{3-p+p^2-p^3}{p^3}.
\end{align}

\begin{theorem}\label{thm:phase}
For every $p\in(0,1)$,
\[
 0<A<B<C<r_0<L<E<U.
\]
The complete target sets $S_{\Ecase},S_{\Lcase}$ for all $r>0$ are
\begin{equation}\label{eq:targetsets}
 \eta^{-1}S_{\Ecase}=
 \begin{cases}
 \{0\},&r<B,\\
 [0,2],&r=B,\\
 \{2\},&B<r<E,\\
 [2,4],&r=E,\\
 \{4\},&r>E,
 \end{cases}
\qquad
 \eta^{-1}S_{\Lcase}=
 \begin{cases}
 \{0\},&r<A,\\
 [0,1],&r=A,\\
 \{1\},&A<r<C,\\
 [1,2],&r=C,\\
 \{2\},&C<r<L,\\
 [2,3],&r=L,\\
 \{3\},&L<r<U,\\
 [3,4],&r=U,\\
 \{4\},&r>U.
 \end{cases}
\end{equation}
In particular, $r_0$ is not an additional first-period phase boundary.
\end{theorem}

For unique targets, $s_{\Lcase}-s_{\Ecase}=\eta$ on $(A,B)$ and
$(L,E)$, while it equals $-\eta$ on $(B,C)$ and $(E,U)$.
Thus demand frequency alone cannot impose a universal temporal ordering
in this family. At equality thresholds, the intervals in
\eqref{eq:targetsets}, rather than just their integer
endpoints, are optimal.

\subsection{Convex clipping}
Rescale to $\eta=h=1$. Let $G(y)=\EE\ell(y-D)$ and
$a=u-pr$. Its slopes on the intervals separated by $0,2$ are $-r,a,1$.
Every objective in \eqref{eq:bellman} is finite, convex and piecewise
linear. For any such $J$, capacity clipping
$T_qJ(x)=\min_{0\le z\le q}J(x+z)$ has right derivative
\begin{equation}\label{eq:clip}
 (T_qJ)'_+(x)=\Phi_qJ'_+(x)
 =\max\{J'_+(x),0\}+\min\{J'_+(x+q),0\}.
\end{equation}
If the minimizer interval of $J$ is $[s_-,s_+]$, the clipped minimizer
interval is $[s_--q,s_+]$. Below that interval optimize at $x+q$;
above it optimize at $x$. This proves \eqref{eq:clip} and handles flat
recourse optima without an arbitrary choice of minimizer.

\subsection{Low cost ratios}\label{sec:low}
Here $a>0$, and $G$ has unique target zero. With $q_3=2$, $V_3$ has
slopes $-r,0,a,1$ at breakpoints $-2,0,2$. The early-loss period-2
objective has slopes
\[
 -(1+p)r\quad(-2,0),\qquad a(2-p)\quad(0,2),\qquad
 2-p+pa\quad(2,4).
\]
As $q_2=0$, $V_2=J_2$, and
\begin{align}\label{eq:elow}
 J'_{1,\Ecase}|_{(0,2)}&=uc-pr(4-2p+p^2)
                      =p(4-2p+p^2)(B-r),\notag\\
 J'_{1,\Ecase}|_{(2,4)}&=c+pa(3-2p)>0.
\end{align}
The objective decreases strictly below zero; hence \eqref{eq:elow}
proves the early-loss alternatives at threshold $B$.

For the late loss, $q_3=1$ and the terminal clipped slopes $-r,0,a,1$
have breakpoints $-1,0,2$. Define
\[
 b_0=u(2-p)-pr(3-p),\qquad R_0=\frac{u(2-p)}{p(3-p)}.
\]
The period-2 slopes on $(0,1),(1,2),(2,3),(3,4)$ are
$b_0,a(2-p),2-p+pa,2-p+pa$. Its target is zero below $R_0$, one
above $R_0$, and the interval $[0,1]$ at equality. Clipping with $q_2=1$
gives
\begin{align}\label{eq:dlow}
 d_0:=J'_{1,\Lcase}|_{(0,1)}&=a+u\max(b_0,0)-p(1+p)r,\notag\\
 d_1:=J'_{1,\Lcase}|_{(1,2)}&=ac+p\min(b_0,0).
\end{align}
Below $R_0$ these are $p(5-3p+p^2)(A-r)$ and $ac>0$;
above $R_0$ they are $u-pr(2+p)$ and $3p(C-r)$.
The identities
\[
 R_0-A=\frac{u(1+p-p^2)}{p(3-p)(5-3p+p^2)}>0,\qquad
 C-R_0=\frac{uc}{3p(3-p)}>0
\]
show that the two current slopes cross zero only at $A$ and $C$.
At $R_0$, $d_0<0<d_1$ and both branch formulas agree.
The derivative below zero is negative and that on $(2,3)$ is positive.
Convexity proves the low-cost late-loss classification globally.

\subsection{High cost ratios}\label{sec:high}
Now $a<0$ and $G$ has unique target 2. For the early loss, $V_3$ has
slopes $-r,a,0,1$ at breakpoints $-2,0,2$. The period-2 target is 2,
and the current slopes are
\begin{equation}\label{eq:ehigh}
 J'_{1,\Ecase}(y)=\begin{cases}
 a(2+p-2p^2)-p(1+p)r<0,&0<y<2,\\
 p^2(1+p)(E-r),&2<y<4,\\
 3-p>0,&4<y<6.
 \end{cases}
\end{equation}
This gives the early-loss transition at $E$, including the whole
minimizer interval $[2,4]$ at equality.

For the late loss, $V_3$ has slopes $-r,a,0,1$ at $-1,1,2$.
The period-2 slopes on the four consecutive unit intervals from 0 to 4 are
\[
 b_0,\quad a(1+p),\quad B_2=2-p^2(1+r),\quad 2-p.
\]
Its target changes from 2 to 3 at $R=2/p^2-1$; at equality it is
$[2,3]$. The current middle slopes are
\begin{equation}\label{eq:dhigh}
 \begin{array}{c|cc}
 & (2,3)&(3,4)\\\hline
 r\le R&2p^2(L-r)&3-3p+p^2\\
 r\ge R&1+p-p^3-p^2(1+p)r&p^3(U-r).
 \end{array}
\end{equation}
The formulas agree at $R$, where the $(2,3)$ slope is $p^2-p-1<0$.
The $(1,2)$ slope is negative: below $R$ it is $a+pb_0<0$, and above
$R$ there is the additional nonpositive term $uB_2$. The $(4,5)$ slope
is positive in both regimes. Since $L<R<U$, the current target crosses
at $L$ and $U$, as asserted.

\subsection{Threshold order and the myopic boundary}
The following positive differences establish every remaining comparison:
\begin{align*}
 B-A&=\frac{u^2c}{p(5-3p+p^2)(4-2p+p^2)},&
 C-B&=\frac{u(3-p+2p^2-p^3)}{3p(4-2p+p^2)},\\
 r_0-C&=u/3,& L-r_0&=\frac{3-3p+p^2}{2p^2},\\
 E-L&=\frac{u(p^2-3p+3)}{2p^2(1+p)},&
 U-E&=\frac{3-p+2p^2-p^3}{p^3(1+p)}.
\end{align*}
Also $R-L=(1+p-p^2)/(2p^2)>0$ and
$U-R=(3-3p+p^2)/p^3>0$.
At $r=r_0$, terminal and recourse optima may be intervals, but the
early current slopes adjacent to 2 are $-(1+p)u<0$ and $c>0$;
the late adjacent slopes are $-pu<0$ and $c>0$. Both current targets
are uniquely 2. This completes the proof of Theorem~\ref{thm:phase}.

\section{Sensitivity and sharpness}\label{sec:sensitivity}
The next bound separates coordinate monotonicity from the false temporal
comparison. Its coordinate statement is prior-known
\cite[Conjecture 2, p.~360]{Journal}; see also \cite{WP,Thesis}.
The journal's Appendix A, p.~368, derives the general bound using
Conjecture 1 and a first-coordinate lemma. The following proof dispenses
with temporal ordering.

\begin{proposition}\label{prop:sensitivity}
In a finite-horizon, zero-order-cost, backlogged inventory system with
linear positive holding/backlog costs, finite first demand moments,
fixed known future capacities and a consistent smallest-minimizer
selector, reducing one capacity by $0\le\eta\le q_i$ gives
\[
 0\le s(q-\eta e_i)-s(q)\le\eta.
\]
Two equal coordinate-loss scenarios relative to one baseline therefore
satisfy $|s(q-\eta e_i)-s(q-\eta e_j)|\le\eta$.
Both signs attain equality in the family of Theorem~\ref{thm:phase},
for every $p\in(0,1)$.
\end{proposition}
\begin{proof}
For a nondecreasing derivative $g$, \eqref{eq:clip} implies
\[
 \Phi_qg(x-\eta)\le\Phi_{q-\eta}g(x)\le\Phi_qg(x).
\]
The clipping operator preserves derivative order and commutes with
translations. Backward expectation preserves this sandwich. Adding the
same convex immediate cost does as well, since its derivative is
nondecreasing. The bounded finite-horizon linear-cost slopes justify
passing right derivatives through expectation. Coercivity supplies
smallest minimizers, obtained as the first point where the right
derivative is nonnegative. Their locations therefore differ by an
amount in $[0,\eta]$. Both loss-scenario targets lie in
$[s(q),s(q)+\eta]$, proving the spread bound. The four strict regions
following Theorem~\ref{thm:phase} attain its two signs.
\end{proof}

The selector condition is substantive at flat optima. The proposition
does not bound the distance between arbitrarily and inconsistently
chosen members of different minimizer intervals.

\section{Exact counterexamples and continuous densities}\label{sec:robust}
For $p=1/2$, the six thresholds are
\[
 A=7/15,\quad B=7/13,\quad C=5/6,\quad
 L=9/2,\quad E=17/3,\quad U=21.
\]
Taking $h=1$, $b=1/2$, $\eta=1$ places $r$ in $(A,B)$, even though
$h>b$. Exact objective values are
\begin{center}
\begin{tabular}{crr}\toprule
$y$ & $J_{1,\Ecase}(y)$ & $J_{1,\Lcase}(y)$\\\midrule
0 & $9/4$ & $17/8$\\
1 & $37/16$ & $33/16$\\
2 & $19/8$ & $5/2$\\\bottomrule
\end{tabular}
\end{center}
The targets are uniquely 0 and 1, while the baseline target is 0.
This contradicts both Conjecture 1 and the separate $h>b$ equal-target
assertion on p.~360 of \cite{Journal}; both also appear in \cite{Thesis}.
For $h=1,b=5$, the targets are uniquely 2 and 3, with optimal objective
values $6$ and $21/4$; the baseline target is 2, with value $9/2$.
Neither example asserts that target ordering determines cost ordering.

\begin{proposition}\label{prop:continuous}
Replace the fair two-point demands in either example by iid
$D_t^\epsilon=2B_t+U_t$, where $U_t$ is uniform on $[0,\epsilon]$
and independent. Both strict target reversals persist for
$\epsilon=10^{-3}$ with an absolutely continuous demand density.
\end{proposition}
\begin{proof}
Let $M=\max(h,b)$. The period loss is $M$-Lipschitz, and a continuation
with $H$ stages has Lipschitz constant at most $HM$. Coupling the two
demand processes and using nonexpansiveness of interval minimization
gives $\sup_y|J_1^\epsilon(y)-J_1^0(y)|\le6M\epsilon$.
If $J_1^0(y)-\min J_1^0\ge\mu|y-y^*|$, every perturbed minimizer
lies within $12M\epsilon/\mu$ of $y^*$. For the low-cost example,
both growth constants are at least $1/16$ and $M=1$, giving radius
$192\epsilon$. For the high-cost example they are at least $1/4$
and $M=5$, giving radius $240\epsilon$. At the stated epsilon the
neighborhoods of the two targets are disjoint.
\end{proof}

\section{A separate loss-location information experiment}\label{sec:information}
Fix either reversal region $(A,B)$ or $(L,E)$. Suppose the early-loss
vector has prior probability $\theta$ and the late-loss vector has
probability $1-\theta$. In one information system the vector is revealed
before period 1; in the other it is revealed after the first demand and
before the period-2 order. All physical capacities, probabilities and
recourse opportunities are otherwise identical. The random loss location
is independent of the entire demand process and of any smoothing noises.
This two-vector prior
induces negative dependence between capacities; it is a separate
common-loss experiment, not the original independent-capacity law.

Let the two targets be $a\eta,(a+1)\eta$, where $a=0$ in the low-cost
reversal and $a=2$ in the high-cost reversal. On the interval between
them scenario regrets have slopes $e>0$ and $-\beta<0$. In the low-cost
region,
\begin{align}\label{eq:infolow}
 e&=hp(4-2p+p^2)(B-r),\notag\\
 \beta&=\begin{cases}
 hp(5-3p+p^2)(r-A),&r\le R_0,\\
 h[pr(2+p)-u],&r\ge R_0;
 \end{cases}
\end{align}
in the high-cost region,
\begin{align}\label{eq:infohigh}
 e&=hp^2(1+p)(E-r),\notag\\
 \beta&=\begin{cases}
 2hp^2(r-L),&r\le R,\\
 h[p^2(1+p)r-1-p+p^3],&r\ge R.
 \end{cases}
\end{align}
Convexity permits restricting the common first decision to this target
interval. The exact expected value of revealing the loss location is
\begin{equation}\label{eq:evpi}
 \mathrm{EVPI}(\theta)=\eta\min\{\theta e,(1-\theta)\beta\}.
\end{equation}
The minimum worst-case regret without a prior is
\begin{equation}\label{eq:minimax}
 \eta\frac{e\beta}{e+\beta},\qquad
 y=\eta\left(a+\frac{\beta}{e+\beta}\right).
\end{equation}
Indeed, a decision $a\eta+d$ incurs regrets $ed$ and
$\beta(\eta-d)$; minimize their maximum or their prior-weighted sum.
At the low-cost numerical example $e=\beta=1/16$, so the fair-prior
EVPI and minimax regret are $1/32$. At the high-cost example each is
$1/8$. These are applications of elementary decision theory within a
precisely matched information experiment, not a new generic information
identity.

\section{Source scope, reproducibility and further questions}
The counterexamples specialize the final journal's Bellman recursion
\cite[Eq.~(4), p.~358]{Journal} to $T=3$, ACI horizon $n=2$,
$\alpha=1$ and $f_4=0$. Its Eq.~(5) and Theorem 2 define the same
unconstrained target and smallest-minimizer selector as used here.
The independent-capacity embedding in Section~\ref{sec:model} respects
the source's periodwise independence assumption. The continuous-density
extension also removes reliance on the interpretation of density notation
for discrete demands.

Version distinctions matter. The final article labels the temporal
comparison conjectural and does not reproduce the earlier working
paper's erroneous stock-translation equality proof. Our journal-specific
conclusions concern Conjecture 1, the adjacent $h>b$ assertion, and the
independent justification of the prior-known sensitivity bound.

Two exact implementations accompany the symbolic proof: a whole-real-line
piecewise-linear solver, and an independently written integer-knot
Bellman audit. Convexity and integer demand/capacity preserve integer
breakpoints, so the latter also certifies flat real minimizer intervals.
An additional continuous scenario-tree LP indexes orders by observed
history and checks the shared first decision in the information experiment.
Its floating-point results are numerical checks, not exact LP proofs.

The six-threshold classification gives sharp conditions within one family.
It leaves open a useful temporal order under wider, explicitly stated
conditions on demand, capacity paths or horizon. It also shows why a
one-coordinate sensitivity bound alone cannot justify truncating distant
capacity information by an asserted temporal ranking.

\begin{thebibliography}{9}
\bibitem{WP}
M. Jak\v{s}i\v{c}, J. C. Fransoo, T. Tan, A. G. de Kok and B. Rusjan,
\emph{Inventory management with advance capacity information},
BETA Working Paper 249, Eindhoven University of Technology, 2008.
\url{https://pure.tue.nl/ws/portalfiles/portal/3296948/691623823560851.pdf}.
\bibitem{Thesis}
M. Jak\v{s}i\v{c}, \emph{Inventory models with uncertain supply},
PhD thesis, Eindhoven University of Technology, 2016, Chapter 2.
\url{https://pure.tue.nl/ws/files/41303873/20161031_Jaksic.pdf}.
\bibitem{Journal}
M. Jak\v{s}i\v{c}, J. C. Fransoo, T. Tan, A. G. de Kok and B. Rusjan,
Inventory management with advance capacity information,
\emph{Naval Research Logistics} \textbf{58} (2011), 355--369.
\href{https://doi.org/10.1002/nav.20450}{doi:10.1002/nav.20450}.
\bibitem{Altug}
M. \c{S}. Altu\u{g} and A. Muharremo\u{g}lu,
Inventory management with advance supply information,
\emph{International Journal of Production Economics} \textbf{129} (2011), 302--313.
\href{https://doi.org/10.1016/j.ijpe.2010.11.005}{doi:10.1016/j.ijpe.2010.11.005}.
\bibitem{Atasoy}
B. Atasoy, R. G\"ull\"u and T. Tan,
Optimal inventory policies with non-stationary supply disruptions and advance
supply information, \emph{Decision Support Systems} \textbf{53} (2012), 269--281.
\href{https://doi.org/10.1016/j.dss.2012.01.005}{doi:10.1016/j.dss.2012.01.005}.
\bibitem{Senga}
J. R. L. Senga, G. A. Chua and D. Chen,
Stochastic inventory systems with regenerative capacity,
\emph{European Journal of Operational Research} \textbf{336} (2027), 152--166;
first available online 7 April 2026.
\href{https://doi.org/10.1016/j.ejor.2026.04.004}{doi:10.1016/j.ejor.2026.04.004}.
\end{thebibliography}
\end{document}
